% ============================================================
%  preamble.tex
%  All packages, font definitions, and custom commands.
%  Edit fonts / colors here; chapter files stay untouched.
% ============================================================

% ── Page geometry ────────────────────────────────────────────
\usepackage[text={4.65in,7.45in}, centering, includefoot]{geometry}

% ── Color ────────────────────────────────────────────────────
\usepackage[table, x11names]{xcolor}

% ── Fonts (XeLaTeX) ──────────────────────────────────────────
\usepackage{fontspec, realscripts}
\usepackage{polyglossia}
  \setdefaultlanguage{english}
  \setotherlanguage{sanskrit}

% Main Latin font (body text — English is the default language)
\setmainfont{Linux Libertine O}

% Named font families — use these in chapter files
\newfontfamily\la  [Script=Devanagari, Scale=.9]{Shobhika}   % Devanagari body
\newfontfamily\regular{Linux Libertine O}
\newfontfamily\en  [Language=English, Script=Latin]{Linux Libertine O}

% --- some other fonts 
% \setmainfont[Script=Devanagari]{Murty Sanskrit}
% \setmainfont[Script=Devanagari]{Siddhanta}
% \setmainfont[Script=Devanagari]{Chandas}
% \setmainfont[Script=Devanagari, Scale=1.1]{Annapurna SIL}
% \setmainfont[Script=Devanagari]{Sanskrit 2003}
% \setmainfont[Script=Devanagari]{Adishila Samskrta}


% Commentary colour fonts — change colors here to restyle all commentaries at once
\newfontfamily\vp  [Script=Devanagari, Scale=1,   Color=purple]{Shobhika-Bold}  % root text / primary
\newfontfamily\vps [Script=Devanagari, Scale=.8,  Color=purple]{Shobhika-Bold}  % small purple
\newfontfamily\vpc [Script=Devanagari, Scale=1,   Color=violet]{Shobhika-Bold}  % secondary
\newfontfamily\qt  [Script=Devanagari, Scale=.8,  Color=violet]{Shobhika-Bold}  % small violet

% Searchable PDF
\XeTeXgenerateactualtext=1

% ── Devanagari page numerals ─────────────────────────────────
\newcommand{\devanagarinumeral}[1]{%
  \devanagaridigits{\number\csname c@#1\endcsname}}

% ── General packages ─────────────────────────────────────────
\usepackage{enumerate}
\usepackage{fancyhdr}
\usepackage{afterpage}
\usepackage{multirow}
\usepackage{multicol}
\usepackage{wrapfig}
\usepackage{caption}
\usepackage{vwcol}
\usepackage{ragged2e}
\usepackage{microtype}
\usepackage{amsmath, amsthm, amsfonts, amssymb}
\usepackage{mathtools}
\usepackage{graphicx}
\usepackage{longtable}
\usepackage[linewidth=1pt]{mdframed}
\usepackage{setspace}
\usepackage{xspace}
\usepackage{array}
\usepackage{emptypage}
\usepackage{hyperref}
  \hypersetup{colorlinks,
    citecolor=black, filecolor=black,
    linkcolor=blue,  urlcolor=black}

% ── Header / footer defaults ─────────────────────────────────
\pagestyle{fancy}
\renewcommand{\headrulewidth}{0pt}  % chapters turn this on as needed

% ── Commentary environment shorthand ─────────────────────────
% Usage:  \begin{commentary}{Commentary Name}  … text …  \end{commentary}
\newenvironment{commentary}[1]{%
  \begin{center}#1\end{center}%
}{}

% ── Convenience: a titled quote block (root-text verses / shlokas) ─────
% Usage:  \roottext{verse text}
\newcommand{\roottext}[1]{%
  \begin{quote}\vpc #1\end{quote}}

% ── Inline Hindi / Devanagari shorthand ──────────────────────
\newcommand{\h}[1]{{\la #1}}

% ── Om symbol ────────────────────────────────────────────────
\newcommand{\om}{{\la ॐ}}
\newcommand{\rt}[1]{\roottext{#1}} % Shortens \roottext to \rt
\newcommand{\omdev}{~॥~\om~॥}       % Auto-inserts the Om signs and dandas

% ── Sarga sub-heading (within a kanda chapter) ───────────────
\newcommand{\sarga}[1]{%
  \vspace{6mm}%
  \begin{center}%
    {\normalsize\textit{#1}}\\%
    \rule{.25\linewidth}{.4pt}%
  \end{center}%
  \vspace{3mm}%
}
